\section{本章小结}
\label{sec:riem-summary}

本章开头提出两个问题：沿流形量出的距离是什么，不同点处的方向怎样比较。第一个问题由黎曼度量回答。度量在每个切空间放一个随位置光滑变化的内积，曲线长度是速率的积分，与参数化无关；内在距离是连接两点的全部分段光滑曲线长度的下确界，它在每个连通分支上是度量，并且不小于嵌入时的环境直线距离。同一光滑流形可以配不同的度量，谈论长度、距离之前须说明所用的度量。

第二个问题需要另一项结构。联络规定向量场沿一个方向怎样求导，它不由光滑结构决定；在黎曼流形上，度量相容与无挠两条要求从中唯一选出Levi-Civita联络，其坐标系数是由度量一阶导数组成的Christoffel符号。沿曲线令协变导数为零得到平行移动$P_\gamma$，它是保持内积的线性同构，但依赖所走的路径。平行移动需要曲线和联络，映射的微分$dF_p$需要点映射，二者是不同的操作。

速度沿自身平行的曲线是测地线。初始位置和速度局部唯一地确定测地线，并光滑地依赖初值，但局部存在不保证能延拓到所有时间；指数映射把从一点出发的测地线汇总成从切空间到流形的映射，在零点附近是微分同胚。固定端点时，三项结论的条件逐步加强：最短且匀速的曲线是测地线，而测地线未必最短；足够短的径向测地线确实最短，而最短测地线未必唯一；完备流形上任意两点之间存在最短测地线，缺少完备性时下确界可能取不到。

曲率用协变导数的次序差扣除Lie括号定义，是逐点的张量，截面曲率在每个切平面上读出一个数：平面与圆柱为零，单位球面为一。它刻画小回路上平行移动的转动，但有限回路还受整体形状影响。内在曲率与嵌入时法向量的变化是两回事，圆与圆柱在环境中弯曲而内在平坦。最后，等距映射是保持切向内积的微分同胚，保持距离、测地线和曲率；等距浸入保持曲线长度，距离却可能缩短。体积、梯度和Laplace--Beltrami算子等依赖度量的微积分运算，在第~\ref{chap:manifold-calculus}章建立。
